
\subsection{Conceptos Relevantes}
La formulación de soluciones parciales y conceptos preliminares se realiza con el objetivo de proyectar posibles direcciones de respuesta sin cerrar anticipadamente la solución final del proyecto. Este apartado permite traducir los hallazgos del análisis en oportunidades de diseño, organizando el problema en dimensiones abordables y estableciendo líneas conceptuales coherentes con el contexto estudiado.

\begin{itemize}
    \item \textbf{Solución parcial 1: Captura ágil del contexto clínico cotidiano.} Diseñar un mecanismo de registro de bajo costo cognitivo que permita al personal asistencial documentar eventos clínicos relevantes en segundos, sin interrumpir el flujo de atención y sin requerir protocolo formal. Este concepto responde al problema de que el contexto cotidiano no se registra porque los sistemas disponibles son demasiado lentos o complejos para el ritmo de trabajo hospitalario.

    \item \textbf{Solución parcial 2: Organización cronológica y trazable del contexto del paciente.} Estructurar los eventos registrados en una bitácora cronológica vinculada al ID de hospitalización del paciente, que permita a cualquier profesional entrante consultar el historial completo del caso sin depender del traspaso verbal. Este concepto aborda el vacío de contexto que experimenta el médico de turno al inicio de cada jornada.

    \item \textbf{Solución parcial 3: Síntesis asistida para el traspaso de turno.} Consolidar automáticamente los eventos del turno en un borrador de evolución que el médico pueda revisar, editar y copiar directamente en la ficha clínica oficial, eliminando la doble documentación y reduciendo el tiempo invertido en el cierre del turno. Este concepto aplica el principio de cero doble trabajo: el sistema genera valor a partir de lo que el profesional ya registró, sin pedirle que lo haga otra vez.
\end{itemize}